\documentclass[11pt]{article}

\usepackage[margin=1in]{geometry}
\usepackage{amsmath,amssymb,amsthm,bm}
\usepackage{graphicx}
\usepackage{xcolor}
\usepackage{float}
\usepackage{listings}
\usepackage{enumitem}
\usepackage{booktabs}
\usepackage{hyperref}
\usepackage{xurl}
\usepackage{fancyhdr}

% ---------- URL and Hyperlink setup ----------
\Urlmuskip=0mu plus 1mu
\hypersetup{
	breaklinks=true,
	colorlinks=true,
	linkcolor=blue,
	urlcolor=blue
}

% ---------- MATLAB code style ----------
\lstdefinestyle{matlabstyle}{
	language=Matlab,
	basicstyle=\ttfamily\small,
	numbers=left,
	numberstyle=\tiny,
	stepnumber=1,
	numbersep=8pt,
	frame=single,
	columns=fullflexible,
	keepspaces=true,
	showstringspaces=false,
	keywordstyle=\bfseries,
	commentstyle=\itshape\color{gray},
}

% ---------- Header setup ----------
\pagestyle{fancy}
\fancyhf{}
\fancyhead[L]{\today}
\fancyhead[C]{ECE 517: Homework 2.2}
\fancyhead[R]{Scott Nguyen}
\renewcommand{\headrulewidth}{0.4pt}
\fancyfoot[C]{\thepage}

\begin{document}
	\setlength{\parindent}{0pt}
	
	Write a complete derivation of the proof that the maximization of margin 
	$d$ is equivalent to minimize the norm of the parameter vector in the maximum margin machine presented in the last set of slides of this module. 
	
	\vspace{1em}
	\textbf{Solution.}
	
	A separating hyperplane can be written as
	\[
	\mathbf{w}^\top \mathbf{x} + b = 0,
	\]
	where $\mathbf{w}$ is the normal vector and $b$ is the bias.
	For linearly separable data $(\mathbf{x}_i, y_i)$ with $y_i \in \{-1, 1\}$, we scale $\mathbf{w}$ and $b$ so that
	\[
	y_i(\mathbf{w}^\top \mathbf{x}_i + b) \ge 1.
	\]
	This means that the two margin planes are
	\[
	\mathbf{w}^\top \mathbf{x} + b = 1
	\quad \text{and} \quad
	\mathbf{w}^\top \mathbf{x} + b = -1.
	\]
	
	Let $\mathbf{x}_0$ be a point on the separating hyperplane:
	\[
	\mathbf{w}^\top \mathbf{x}_0 + b = 0.
	\]
	A point on the upper margin can be written as $\mathbf{x}_1 = \mathbf{x}_0 + \rho \mathbf{w}$.
	Substituting into $\mathbf{w}^\top \mathbf{x}_1 + b = 1$ gives
	\[
	\mathbf{w}^\top (\mathbf{x}_0 + \rho \mathbf{w}) + b = 1
	\quad \Rightarrow \quad
	\rho \|\mathbf{w}\|^2 = 1
	\quad \Rightarrow \quad
	\rho = \frac{1}{\|\mathbf{w}\|^2}.
	\]
	The perpendicular distance from the hyperplane to this margin is
	\[
	d = \|\mathbf{x}_1 - \mathbf{x}_0\|
	= \rho \|\mathbf{w}\|
	= \frac{1}{\|\mathbf{w}\|}.
	\]
	
	Since $d = 1/\|\mathbf{w}\|$, maximizing the margin is the same as minimizing $\|\mathbf{w}\|$.
	So the maximum–margin problem can be written as
	\[
	\min_{\mathbf{w}, b} \|\mathbf{w}\|
	\quad \text{subject to} \quad
	y_i(\mathbf{w}^\top \mathbf{x}_i + b) \ge 1.
	\]
	
	\medskip
	Therefore, maximizing the margin $d$ is equivalent to minimizing the norm of the parameter vector $\mathbf{w}$ in the maximum–margin machine.
	
\end{document}